\chapter{Literature Review}

\section{Overview}
This chapter reviews the research areas that inform \projectname: image-based virtual try-on, multiview try-on and view synthesis, generative models for conditioned image synthesis, 3D reconstruction, human body understanding and measurement estimation, and retrieval-augmented discovery for fashion e-commerce. Rather than surveying each topic in isolation, the discussion traces how methods in each area have evolved, identifies their respective limitations, and highlights how their integration motivates a unified fashion platform. The progression from single-view garment transfer toward multiview-consistent, measurement-aware, and retrieval-informed shopping experiences forms the central thread of this review.

\section{Image-Based Virtual Try-On}
Image-based virtual try-on aims to synthesize a realistic image of a person wearing a target garment, given only a person photograph and a garment image as input. Early approaches decompose this problem into two stages: geometric alignment of the garment to the person's body, followed by pixel-level refinement to blend the warped garment into the scene.

VITON introduced the foundational two-stage pipeline, combining a coarse result from a clothing-agnostic representation with a refinement network that sharpens textures and corrects boundary artifacts \cite{han2018viton}. While VITON demonstrated the viability of image-conditioned garment transfer, it suffered from noticeable misalignment under non-trivial poses and produced blurred textures in complex garment regions. CP-VTON improved upon this by introducing a geometric matching module based on thin-plate spline (TPS) transformation, which better aligns the garment to the target pose before synthesis \cite{wang2018cpvton}. The explicit geometric warping step reduced spatial misalignment, but the approach remained sensitive to large pose differences and struggled with fine-grained texture preservation.

Subsequent methods shifted toward higher-resolution outputs and more robust warping strategies. HR-VITON addressed the resolution bottleneck by incorporating a misalignment-aware segmentation module and a conditional generator that operates at higher spatial dimensions \cite{lee2022hrviton}. This approach significantly improved visual quality at 1024$\times$768 resolution, though it introduced higher computational cost and required carefully curated training data at the target resolution.

The emergence of diffusion-based generative models marked a further shift in try-on quality. StableVITON adapted latent diffusion models for virtual try-on by introducing zero-cross-attention blocks that inject garment features into the denoising process without requiring explicit warping \cite{kim2024stableviton}. This eliminated the need for separate geometric transformation modules and improved texture fidelity, particularly for patterned and textured garments. However, the approach inherited the computational cost of diffusion sampling, making real-time deployment challenging.

IDM-VTON extended the diffusion-based paradigm by combining high-level garment semantics with low-level texture features through a dual-encoder architecture \cite{choi2024idmvton}. The model uses a TrueNet component for identity preservation and achieves state-of-the-art results on standard benchmarks. Nevertheless, like other diffusion-based methods, inference latency remains a practical concern for interactive applications.

OOTDiffusion proposed an alternative approach by formulating try-on as outfitting fusion within a pretrained latent diffusion framework \cite{xu2024ootdiffusion}. Rather than conditioning the diffusion process on warped garment features, OOTDiffusion fuses garment appearance through outfitting dropout and attention-based feature injection. This design preserves garment details more faithfully than warp-based methods and supports both half-body and full-body try-on scenarios. However, OOTDiffusion operates in single-view mode and does not address multiview consistency, limiting its applicability for 360-degree visualization.

Across these methods, a clear progression emerges: from explicit geometric warping (VITON, CP-VTON) to learned alignment at higher resolutions (HR-VITON), and finally to diffusion-based synthesis that bypasses warping entirely (StableVITON, IDM-VTON, OOTDiffusion). Each generation improves texture fidelity and pose robustness, but introduces new tradeoffs in computational cost, training data requirements, and deployment complexity. The challenge of producing consistent results across multiple viewpoints remains largely unaddressed by single-view methods, motivating the multiview extensions discussed in the following section.

\section{Multiview Virtual Try-On}
Single-view try-on outputs, while increasingly realistic, provide only a partial visualization of how a garment looks on a person. In real shopping scenarios, consumers benefit from seeing the garment from multiple angles, which requires multiview consistency—the garment must appear coherent in color, texture, and fit across different viewpoints.

VTON 360 formulates 3D try-on as a multiview extension of 2D pipelines \cite{he2025vton360}. The method introduces view-consistent conditioning to preserve garment identity across viewpoints, using multiview priors and cross-view attention mechanisms to enforce coherence between independently generated views. This approach demonstrates that multiview conditioning and geometric priors can significantly improve consistency compared to naive per-view synthesis, where each viewpoint is generated independently and often exhibits inconsistent garment appearance, color shifts, or structural discontinuities.

However, multiview try-on methods face several challenges that remain active areas of research. First, maintaining fine-grained texture consistency across large angular differences (for example, between front and back views) is difficult because different viewpoints expose different garment regions with minimal overlap. Second, occlusion handling becomes more complex in side and back views, where the garment may be partially hidden by the body or other clothing. Third, the computational cost scales with the number of views, as each view typically requires a separate forward pass through the generation model. These challenges motivate the use of explicit 3D representations for view synthesis, discussed in Section~2.5, as a complementary approach to multiview-conditioned 2D generation.

\section{Generative Models for Try-On}
The quality improvements observed in recent try-on systems are closely tied to advances in conditional generative modeling, particularly diffusion models. Understanding these foundations is essential for contextualizing the architectural choices in modern try-on pipelines.

Denoising diffusion probabilistic models (DDPMs) introduced a principled framework for iterative image generation through a forward noising process and a learned reverse denoising process \cite{ho2020ddpm}. DDPMs produce high-fidelity samples and offer stable training dynamics compared to earlier generative adversarial networks (GANs), but their iterative sampling requires hundreds of denoising steps, resulting in slow inference.

Latent diffusion models (LDMs) addressed this computational bottleneck by performing the diffusion process in a compressed latent space rather than directly in pixel space \cite{rombach2022ldm}. By encoding images into a lower-dimensional latent representation using a pretrained autoencoder, LDMs reduce memory and computation requirements while preserving perceptual quality. The introduction of cross-attention conditioning in LDMs enables flexible integration of external signals such as text prompts, segmentation masks, or reference images, making LDMs particularly well-suited for conditioned synthesis tasks including virtual try-on.

These advances have directly influenced the try-on literature. StableVITON, IDM-VTON, and OOTDiffusion all build upon the LDM framework, leveraging its conditioning mechanisms to inject garment information into the generation process. The key design question across these systems is how to encode garment appearance—whether through warped feature maps, cross-attention injection, or outfitting fusion—such that the generated output faithfully reproduces garment details while respecting the person's pose and body shape. This interplay between generative model architecture and garment conditioning strategy remains a central research challenge, with each approach offering different tradeoffs between fidelity, consistency, and inference speed.

\section{360 View Synthesis and 3D Reconstruction}
While multiview try-on generates a discrete set of viewpoint images, full 360-degree visualization requires continuous novel view synthesis or explicit 3D reconstruction. This capability enables interactive product previews where users can freely rotate the visualized result, providing a richer understanding of garment fit and appearance.

Neural Radiance Fields (NeRF) represent a scene as a continuous volumetric function that maps 3D coordinates and viewing directions to color and density values \cite{mildenhall2020nerf}. By training on a sparse set of posed images, NeRF can synthesize photorealistic novel views through volumetric rendering. NeRF demonstrated remarkable visual quality for view synthesis, but its reliance on per-scene optimization and expensive volumetric ray marching makes it impractical for real-time applications. Training a single NeRF scene can take hours, and rendering a single frame requires evaluating the network at hundreds of sample points per ray.

3D Gaussian Splatting (3DGS) addressed these limitations by representing scenes as collections of 3D Gaussian primitives that are rasterized using a differentiable splatting operation \cite{kerbl2023gaussians}. This explicit point-based representation enables real-time rendering at high resolution while maintaining visual quality comparable to NeRF. Gaussian splatting also supports faster optimization, with training times measured in minutes rather than hours for typical scenes. The Nerfstudio framework provides modular implementations of both NeRF and Gaussian splatting variants, including Splatfacto, which offers a production-ready Gaussian splatting pipeline with configurable densification, pruning, and appearance modeling strategies \cite{tancik2023nerfstudio}.

For fashion applications, the choice between NeRF-based and splatting-based methods involves clear tradeoffs. NeRF provides smoother interpolation in regions with sparse observations but cannot achieve interactive frame rates. Gaussian splatting achieves real-time rendering and faster training but may produce artifacts in under-observed regions where Gaussian primitives are insufficiently constrained. For an interactive e-commerce platform where users expect responsive 360-degree previews, the real-time rendering capability of Gaussian splatting is a decisive advantage, even at the cost of occasional visual artifacts in challenging viewpoints.

\section{Human Body Understanding and Preprocessing}
Virtual try-on and body measurement pipelines require accurate understanding of human body structure as a preprocessing step. This involves detecting and segmenting the person in the input image, estimating body pose, and extracting dense correspondences between the image and a canonical body surface. The quality of these preprocessing outputs directly affects the accuracy of downstream tasks.

DensePose extends the problem of human pose estimation from sparse keypoints to dense surface correspondences \cite{guler2018densepose}. Given an input image, DensePose predicts a mapping from every pixel belonging to a person to a UV coordinate on the SMPL body surface \cite{loper2015smpl}. This dense correspondence map provides the geometric foundation for garment warping in try-on pipelines, as it enables precise spatial alignment between the input image and a canonical body representation. DensePose predictions are typically produced using architectures built on top of the Detectron2 framework, which provides modular object detection and segmentation components \cite{wu2019detectron2}.

Person detection, which precedes DensePose and parsing, is commonly handled by general-purpose object detectors. The YOLO family of detectors, including YOLOv8, provides fast and accurate bounding box detection suitable for identifying persons in unconstrained images \cite{redmon2016yolo,jocher2023yolov8}. In preprocessing pipelines, person detection serves as the initial step that crops and isolates the relevant region before applying more expensive per-pixel analysis such as DensePose and human parsing.

Human parsing complements DensePose by providing semantic segmentation of clothing regions, labeling each pixel as belonging to categories such as upper body, lower body, dress, hair, or skin. Together, DensePose correspondences and parsing maps provide the conditioning inputs that modern try-on models require to accurately place garments on the person's body.

The accuracy and robustness of these preprocessing components are critical. Errors in person detection lead to incorrect crops, DensePose failures produce garbled surface mappings, and parsing mistakes cause garments to bleed into skin or hair regions. In practice, building a reliable preprocessing pipeline requires careful orchestration of multiple models and handling of edge cases such as occluded bodies, unusual poses, and low-resolution inputs.

\section{Body Measurement Estimation}
Accurate size recommendation in fashion e-commerce requires estimating body measurements from visual input. Traditional approaches rely on manual measurement or depth sensors, but image-based methods offer a more accessible alternative that requires only a standard photograph.

Parametric human body models provide the geometric foundation for measurement estimation. SMPL represents the human body as a differentiable function of shape and pose parameters, producing a 3D mesh from a compact parameter vector \cite{loper2015smpl}. Shape parameters control body proportions such as height, weight distribution, and limb lengths, while pose parameters control joint rotations. By fitting SMPL parameters to visual observations, systems can recover a 3D body mesh from which anthropometric measurements can be directly computed.

SHAPY advanced this approach by training a regression network that predicts SMPL body shape parameters directly from a single image, using linguistic shape attributes as additional supervision \cite{choutas2022shapy}. Unlike earlier fitting-based methods that require iterative optimization, SHAPY produces shape estimates in a single forward pass, enabling faster inference. The model was trained on a dataset annotated with natural language descriptions of body shape, which provides a richer supervisory signal than mesh-only losses. SHAPY demonstrates that accurate body shape estimation—and by extension, measurement extraction—is feasible from monocular images without requiring depth sensors or multiple viewpoints.

However, image-based measurement estimation faces inherent limitations. Clothing, pose variation, and camera perspective introduce systematic biases that affect circumference and width estimates. Loose clothing obscures the true body contour, leading to overestimation of measurements. Camera distance and lens distortion affect apparent proportions, requiring calibration with a known reference measurement such as the person's height. These challenges motivate hybrid approaches that combine model-predicted shape with user-provided reference values to correct for systematic errors, an approach adopted in the system described in this thesis.

\section{Retrieval-Augmented Discovery for Fashion E-Commerce}
Product discovery in fashion e-commerce requires systems that understand natural language queries, match them against large product catalogs, and return relevant results with explanations. Traditional keyword-based search fails to capture the semantic richness of fashion queries, which often involve abstract concepts such as style, occasion, and aesthetic preference rather than exact product names.

Retrieval-augmented generation (RAG) provides a framework for combining retrieval and generation in a single pipeline \cite{lewis2020rag}. In the RAG paradigm, a retriever first identifies relevant documents from an external knowledge base, and a generator then produces a response conditioned on both the user query and the retrieved context. This architecture grounds generated responses in factual catalog data, reducing hallucination and ensuring that recommendations correspond to real, available products.

The retrieval component in fashion applications typically relies on dense embedding models. Sentence-BERT produces fixed-dimensional sentence embeddings through a Siamese network architecture, enabling efficient similarity computation via cosine distance \cite{reimers2019sbert}. When applied to product catalogs, each product is embedded based on its textual description (name, material, color, style tags), and user queries are embedded in the same space. Retrieval then becomes a nearest-neighbor search in embedding space, which scales to large catalogs using vector databases or approximate nearest-neighbor indices.

Fashion-specific retrieval introduces additional challenges beyond general semantic search. Product descriptions are often short and formulaic, requiring the embedding model to capture subtle distinctions between similar items. Color matching requires fuzzy expansion, as users may request ``blue'' while the catalog contains entries labeled ``navy,'' ``cobalt,'' or ``sky blue.'' Size and availability constraints introduce hard filters that must be applied before or alongside semantic ranking. Budget and occasion requirements add further dimensions that pure embedding similarity cannot capture.

These considerations motivate hybrid retrieval architectures that combine structured metadata filtering with semantic ranking. A two-stage pipeline—first applying hard filters on attributes such as stock, gender, price range, and size availability, then ranking the filtered candidates by semantic similarity—balances precision and recall while remaining computationally efficient. Progressive filter relaxation, where constraints are loosened if too few candidates survive the initial filter, ensures that the system degrades gracefully rather than returning empty results.

Intent classification adds another layer of intelligence to the retrieval pipeline. Not every user message is a product search; greetings, follow-up questions, and non-fashion queries require different handling. Classifying user intent before invoking the retrieval pipeline avoids unnecessary computation and enables the system to respond appropriately to conversational inputs. This classification step, combined with structured intent extraction that identifies specific attributes such as garment type, color, and occasion from the user's message, transforms an unstructured natural language query into a structured retrieval request.

The integration of retrieval, intent understanding, and response generation into a coherent pipeline connects the RAG literature to the broader goal of intelligent shopping assistance. Rather than treating product search as an isolated information retrieval problem, the system must orchestrate multiple components—classification, extraction, retrieval, ranking, and generation—into a responsive conversational experience. This systems integration perspective, rather than any single algorithmic contribution, distinguishes applied fashion RAG systems from the foundational NLP research on which they build.

\section{Summary}
The literature reviewed in this chapter spans several distinct but interconnected research areas. Image-based virtual try-on has progressed from explicit geometric warping to diffusion-based synthesis, steadily improving texture fidelity and pose robustness while introducing new tradeoffs in computational cost. Multiview extensions address the consistency gap left by single-view methods but remain computationally expensive and sensitive to large viewpoint changes. 3D reconstruction techniques, particularly Gaussian splatting, enable real-time 360-degree visualization that complements multiview try-on outputs. Human body understanding through dense pose estimation and body measurement regression provides the geometric and anthropometric foundations that try-on and sizing systems require. Finally, retrieval-augmented discovery combines semantic search with structured filtering and conversational intent understanding to support intelligent product exploration.

No single existing system integrates all of these capabilities into a unified platform. Individual research contributions advance specific components—better try-on quality, faster 3D rendering, more accurate measurements, smarter retrieval—but they are typically evaluated in isolation. The contribution of \projectname\ lies in combining these components into an integrated architecture that supports the full shopping workflow: discovery, visualization, measurement, and purchase. The design decisions and implementation details of this integration are presented in the subsequent chapters.
